% =============================================================================
% Executive Summary
% =============================================================================
\section*{Executive Summary}\label{sec:executive-summary}
\addcontentsline{toc}{section}{Executive Summary}

This paper identifies the dominant structural driver of cross-country
differentiation in the International Sovereignty Index (ISI) applied to
the EU-27.  The finding is concise and falsifiable:

\begin{tcolorbox}[
  colback=ISILight,
  colframe=ISIBlue,
  fonttitle=\bfseries\sffamily,
  title={Core finding}
]
Two of the six {\ISI} axes---defence and logistics---account for
69.4\% of total cross-country composite variance.  The remaining four
axes (financial, energy, technology, critical inputs) contribute
30.6\% collectively.  In all 27~countries, either defence or logistics
is the highest-scoring axis; no country's profile is dominated by any
other axis.
\end{tcolorbox}

\medskip
\noindent
The principal results are:

\begin{itemize}[leftmargin=2em]

\item \textbf{Variance dominance.}
  A covariance-based variance decomposition attributes 35.1\% of
  composite variance to the defence axis and 34.3\% to the logistics
  axis.  The own-variance component accounts for 64.3\% of total
  composite variance; cross-covariances contribute 35.7\%.
  Financial concentration accounts for 1.0\%.
  See \cref{tab:variance-decomp} and \cref{fig:variance-bar}.

\item \textbf{Universal axis dominance.}
  Defence is the highest-scoring axis in 16 of 27~countries;
  logistics is highest in the remaining~11.  No country is dominated
  by financial, energy, technology, or critical inputs.  Among
  defence-dominant countries, the mean gap between the highest and
  second-highest axis is 0.184 on the unit interval.
  See \cref{tab:dominance-counts} and \cref{fig:dominance-share}.

\item \textbf{Rank disruption under omission.}
  Removing the defence axis from the composite produces the largest
  rank disruption: Spearman~$\rho = 0.833$, Kendall~$\tau = 0.675$.
  Removing financial or energy leaves the ranking nearly unchanged
  ($\rho > 0.99$).
  See \cref{tab:loo-disruption} and \cref{fig:loo-rho}.

\item \textbf{Weight robustness.}
  Dirichlet weight perturbation (10\,000 draws, $\alpha_j = 10$)
  yields a mean rank-order correlation of $\rho = 0.980$.  The
  Netherlands exhibits the highest rank volatility ($\sigma = 3.08$).
  See \cref{tab:dirichlet-top}.

\item \textbf{Correlation structure.}
  Only 2 of 15~pairwise axis correlations are significant after
  Bonferroni correction.  The financial axis is orthogonal to all
  others ($|t| < 1.15$).

\end{itemize}

\medskip
\noindent
The implication for policy is direct: any reform that alters the
defence or logistics axis will propagate to the composite ranking;
changes to the remaining four axes are, for ranking purposes,
second-order.  For measurement, the result implies that the index's
discriminatory power rests on two domains where data are structurally
sparse (arms transfers) or operationally opaque (maritime and freight
logistics).
